\documentclass[11pt]{article}
\usepackage{icmlko}
\begin{document}

\title{학습 행을 버려서 얻었다 --- 만화 절반 채택, 그리고 ``차가 정확히 $0$''이 잡아낸 배선 버그}
\author{월드모델 실험실}
\date{노트 509 · 2026년 8월 3일}
\maketitle

\begin{abstract}
노트 508 이 도메인별 행당 값어치를 재다가 만화 하나만 부호가 반대인 것을
봤다 --- 학습 행 $2{,}823$ 개(전체 학습의 $15\%$)를 빼면 판이 \emph{오른다}.
관문에 태웠다.

\begin{center}
\begin{tabular}{lrrl}
\hline
자 & 차 & 부호 & $95\%$ 구간 \\
\hline
판 & $+0.0021$ & $30/40$ & $[+0.0013, +0.0030]$ \\
날짜 통제 판 & $+0.0035$ & $34/40$ & $[+0.0024, +0.0046]$ \\
무리 안 판 & $+0.0022$ & $31/40$ & $[+0.0011, +0.0032]$ \\
\textbf{KR 만화} & $\mathbf{+0.0274}$ & $38/40$ & $[+0.0227, +0.0321]$ \\
\textbf{비게임 앱} & $\mathbf{+0.0082}$ & $33/40$ & $[+0.0049, +0.0115]$ \\
\hline
\end{tabular}
\end{center}

\textbf{자 다섯이 전부 양수이고 전부 $0$ 밖이다.} 그리고 노트 507 이
어제 만든 조항(판 차 $+0.002$ 아래면 씨앗 $40$ 까지)을 지켜 씨앗 $20 \to
40$ 으로 갔더니 \textbf{더 단단해졌다} --- 세계애니 \texttt{grp} 이
같은 검사에서 무너진 것과 정반대다.

\textbf{배포 모양(6-앙상블)에서 더 크다} --- $+0.0033$($6/6$,
$[+0.0026, +0.0040]$) · 날짜 통제 $+0.0052$($6/6$). 판 $0.4852 \to
\mathbf{0.4885}$. \textbf{채택한다} --- \texttt{TRAINCAP = \{"만화": 0.5\}}.
\end{abstract}

\begin{figure}[t]
\centering
\includegraphics[width=\textwidth]{figs/f1.pdf}
\caption{\textbf{가} --- 씨앗을 늘리며. \textbf{나} --- 자 다섯.
\textbf{다} --- 만화의 몫.}
\end{figure}

\section{왜 만화인가}

\begin{itemize}
\item \textbf{시끄러워서가 아니다} --- 만화의 축$\leftrightarrow$라벨
$|rho|$ 평균이 $0.197$ 로 남($0.195$)과 같다. 부호가 갈리는 축도 둘뿐이고
작다(\texttt{entry\_friction} $-0.158$ · \texttt{trend\_momentum}
$-0.062$).
\item \textbf{과대표집이 유일한 설명도 아니다} --- 과대표집 배율과 행당
값어치의 스피어만이 $0.143$ 이다.
\item \textbf{그러나 만화는 극단이다} --- 학습의 $\mathbf{30.5\%}$ 인데
채점의 $\mathbf{7.5\%}$ 다(비 $4.05$, 다음인 세계애니 $1.63$ 의 $2.5$ 배).
그리고 노트 505 는 만화 챔피언이 자기 천장의 $\mathbf{2.00}$ 배라고 했다.
\textbf{제일 많이 받고 유일하게 안 준다.}
\end{itemize}

\textbf{집 밖 이득이 판 이득의 $13$ 배}인 것이 그 그림과 맞는다 --- 만화
행이 공유 축을 자기 쪽으로 끌어당기고 있었고, 그 왜곡은 집안 채점보다
\emph{집 밖}에서 훨씬 아프다.

\section{``차가 정확히 $+0.0000$''}

첫 앙상블 확인에서 세 블록이 모두 \textbf{차 $+0.0000$} 을 냈다. 효과가
없는 것이 아니라 \textbf{패치가 안 걸린 것}이었다 ---
\texttt{str.replace(old, new, 1)} 이 \emph{첫 번째} 일치를 바꾸는데, 부모
클래스 \texttt{Boost.fit} 의 첫 세 줄이 자식 \texttt{BagBoost.fit} 과
똑같아서 부모가 맞았다. 챔피언은 자식을 쓴다.

\begin{quote}
\textbf{조항} --- \textbf{두 설정의 차가 \emph{정확히} $0$ 이면 효과가
없다고 읽지 말고 배선을 먼저 의심한다.} 잡음이 있는 자에서 정확한 $0$ 은
자연스럽지 않다 --- 그것은 두 실행이 \emph{같은 코드}였다는 표시다.
\end{quote}

부수 조항 --- \textbf{상속받은 메서드를 문자열로 패치할 때는 클래스 범위
안에서 찾는다.} 부모와 자식이 같은 첫 줄을 가지면 \texttt{count=1} 이
부모를 맞춘다.

\section{규약 둘이 연달아 일했다}

\begin{itemize}
\item \textbf{노트 507 의 문턱}($+0.002$ 아래면 씨앗 $40$) --- 이 후보가
$+0.0020$ 으로 문턱에 \emph{정확히} 걸렸고, 그래서 $40$ 까지 갔고,
통과했다. 어제 그 조항을 만든 이유(세계애니 \texttt{grp} 이 $20$ 에서
살아 있다가 $40$ 에서 죽음)와 정반대 결과다. \textbf{문턱은 죽이기도
살리기도 한다.}
\item \textbf{배포 모양으로 다시 재기} --- 관문은 단일 모형인데 배포는
앙상블이다. 이번엔 앙상블에서 \emph{더} 컸지만, 안 그랬으면 채택을
되돌렸을 것이다. \textbf{새 조항: 채택은 배포 모양에서 한 번 더 잰다.}
\end{itemize}

\section*{남은 것}

\begin{itemize}
\item \textbf{절반이 최적인가} --- $25\% \cdot 75\%$ 는 안 재 봤다.
\item \textbf{만화만인가} --- 세계애니도 과대표집 $1.63$ 배다. 같은 자로
태울 수 있다.
\item \textbf{원리적인 판} --- 학습 행을 유보 무게에 맞춰 \emph{가중}하는
쪽이 도메인마다 문턱을 정하는 것보다 낫다. \texttt{sample\_weight} 로
할 수 있다.
\end{itemize}

\end{document}
